\scanpage{98}
% archive foliation: 97
% A faint upper-right mark resembles 1; not securely distinguishable from later foliation.
Cartes \inserted{\uline{totalement}} épinglées. etc.

I) Appelons pré-carte un système $(S,A,F\deleted{\unsure{\}}};p,q,r)$ où $S,A,F$ sont des ensembles, $p,q,r$ des applications
\deleted{de $S,A,$}
\[
S\xrightarrow{p}\mathbb N^{*}\qquad A\xrightarrow{q}\mathbb N^{*}\qquad F\xrightarrow{r}\mathbb N^{*}
\]
et $q$ ne prenant que les valeurs $1$ et $2$.

\unsure{Ils|Il} forment une catégorie pour les \deleted{iso.}\inserted{morphismes} $\varphi$ de $(S',A',F';p',q',r')\longrightarrow(S,A,F,p,q,r)$
\deleted{(\unsure{cartes}) \unsure{triples} d'}
\deleted{À toute carte combinatoire $C$, on associe \unsure{une pré-carte}}
\inserted{applications}
\deleted{$\varphi_s:S'\;\unsure{-|\longrightarrow}\;S$}
\[
\varphi_s:S'\longrightarrow S,\quad\varphi_a:A'\;\unsure{-|\longrightarrow}\;A,\quad\varphi_f:F'\;\unsure{-|\longrightarrow}\;F
\]
satisfaisant
\[
\left\{\begin{aligned}
&\deleted{p(\varphi_s(s'))\mid p'(\unsure{s'})}\quad\inserted{p'(s')=\nu(\varphi,s)\,p(s)}
&&\forall s'\in S',\ s=\varphi_s(s')\\
&\deleted{q(\varphi_a(a'))\mid q'(\unsure{a'})}\quad\inserted{q'(a')=\nu(\varphi,a)\,q(a)}
&&\forall a'\in A',\ a=\varphi_a(a')\\
&\deleted{r(\varphi_f(f'))\mid r'(\unsure{f'})}\quad\inserted{r'(f')=\nu(\varphi,\unsure{f|r})\,r(f)}
&&\forall f'\in F',\ f=\varphi_f(f')
\end{aligned}\right.
\]
où $\nu(\varphi,s)=\operatorname{Card}(\varphi_s^{-1}(s))$, $\nu(\varphi,a)=\text{---}$, $\nu(\varphi,f)=\text{---}$.

On a un foncteur évident de la catégorie \inserted{$(\mathrm{Cartcomb})$} des cartes comb. (avec bords, \inserted{\unsure{facettes pendantes}} \unsure{plis}) dans celle des précartes, $(\mathrm{Précart})$,
\[
C\longmapsto\operatorname{Fac}(C).
\]
\deleted{Appelons épinglage d'}
Ce foncteur est conservatif\margin{$\operatorname{Fac}(u)$ iso $\Rightarrow u$ iso}.
\margin{\unsure{D'où} \uline{\unsure{fidèle}}\quad \unsure{ce foncteur est} \unsure{fidèle} pour iso}

Soit $\mathcal C_0=(S,A,F,p,q,r)$ une pré-carte.
On appelle carte \inserted{comb.} $\mathcal C_0$-épinglée la donnée d'une carte $C$ et d'un iso.
\[
\deleted{\unsure{\operatorname{Fac}(C)}\longrightarrow\unsure{C}}
\qquad\mathcal C_0\xrightarrow[\sim]{u}\operatorname{Fac}(C)
\]
\deleted{\unsure{C'est une}} Les morphismes sont les morphismes de cartes compatibles avec l'épinglage $\mathcal C_0$, bien sûr des isomo. de cartes, et il se
